% ===========================================================================
% information_limits.tex
% SCX Information-Theoretic Limits
% Kolmogorov Complexity of Audit Claims | Optimal Hypothesis Set Length Lower Bound | SCX Generalization of Quantum Fisher Information
% Honest-Punch Edition
% ===========================================================================
\documentclass[12pt,a4paper]{article}

% ---- Packages ----
\usepackage{amsmath,amssymb,amsthm}
\usepackage{mathtools}
\usepackage{booktabs}
\usepackage{geometry}
\usepackage{hyperref}
\usepackage{enumitem}
\usepackage{tikz}
\usepackage{tcolorbox}
\tcbuselibrary{skins,breakable}

\geometry{left=2.5cm,right=2.5cm,top=2.5cm,bottom=2.5cm}

% ---- Theorem environments ----
\theoremstyle{plain}
\newtheorem{theorem}{Theorem}[section]
\newtheorem{lemma}[theorem]{Lemma}
\newtheorem{corollary}[theorem]{Corollary}
\newtheorem{proposition}[theorem]{Proposition}

\theoremstyle{definition}
\newtheorem{definition}[theorem]{Definition}
\newtheorem{assumption}[theorem]{Assumption}
\newtheorem{remark}[theorem]{Remark}

% ---- Strictness markers ----
\newcommand{\strictness}[1]{{\small\textsf[#1]}}
\newcommand{\rigorous}{\strictness{Rigorous: verifiable line by line}}
\newcommand{\sketch}{\strictness{Sketch: direction correct, details TBD}}
\newcommand{\conjecture}{\strictness{Conjecture: core steps unproven}}
\newcommand{\heuristic}{\strictness{Heuristic: intuition-driven}}

% ---- Honest commentary boxes ----
\newtcolorbox{honestbox}[1][]{
  colback=red!5!white, colframe=red!75!black, fonttitle=\bfseries,
  title={Honest Punch}, breakable, #1
}

\newtcolorbox{limitbox}[1][]{
  colback=violet!3!white, colframe=violet!60!black, fonttitle=\bfseries,
  title={Limit Test}, breakable, #1
}

% ---- Operators ----
\DeclareMathOperator{\Kolm}{K}
\DeclareMathOperator{\Ent}{H}
\DeclareMathOperator{\MI}{I}
\DeclareMathOperator{\KL}{D_{KL}}
\DeclareMathOperator{\Var}{Var}
\DeclareMathOperator{\Cov}{Cov}
\DeclareMathOperator{\Tr}{Tr}
\DeclareMathOperator{\rank}{rank}
\newcommand{\complexity}[1]{\mathcal{K}(#1)}
\newcommand{\FisherClassical}{\mathcal{I}}
\newcommand{\FisherQuantum}{\mathcal{F}}

\title{\textbf{SCX Information-Theoretic Limits}\\[0.3em]
       \large Kolmogorov Complexity of Audit Claims, Optimal Hypothesis Set Length Lower Bound,\\SCX Generalization of Quantum Fisher Information}
\author{SCX Theory Group}
\date{\today}

\begin{document}
\maketitle

\begin{abstract}
This paper explores the intersection of the SCX framework with \textbf{three deep information-theoretic limits}:

(1) \textbf{Kolmogorov Complexity:} The effective testability of any SCX audit claim is bounded by the algorithmic complexity of the claim itself --- ``an audit assertion that cannot be generated by a short program cannot be tested by short evidence.''

(2) \textbf{Optimal Hypothesis Set Length Lower Bound:} SCX's minimal sufficient statistic is jointly determined by the structural depth of the CEC chain and the resolution of the audit target --- there exists an unconditional information-theoretic lower bound of $\Omega(\text{CDepth} \cdot \log|\mathcal{A}|)$.

(3) \textbf{SCX Generalization of Quantum Fisher Information:} In a quantum SCX extension (if CEC states are quantum states), the quantum-mechanical limit of audit sensitivity is given by the generalized quantum Fisher information matrix.

\textbf{Honest Punch:} Among the three limits, (1) is a restatement of classical AIT in SCX terms (not original), (2) has a rigorous proof but the lower bound is overly loose (by a $\log$ factor), and (3) is currently only a \textbf{formal analogy} --- the existence of quantum SCX itself has not been defined.
\end{abstract}

\tableofcontents

% ===========================================================================
\section{Kolmogorov Complexity of Audit Claims}
\label{sec:kolmogorov}
% ===========================================================================

\begin{definition}[Kolmogorov Complexity]
\label{def:kolmogorov}
The (prefix) Kolmogorov complexity $\Kolm(x)$ of a string $x$ is defined as the length of the shortest prefix Turing machine program that outputs $x$ and halts:
\begin{equation}
\Kolm(x) = \min\{|p| : U(p) = x\},
\end{equation}
where $U$ is a universal prefix Turing machine. $\Kolm(x)$ is well-defined up to $O(1)$ independent of the choice of $U$ (invariance theorem).
\end{definition}

\begin{theorem}[Complexity Lower Bound for Audit Claims --- SCX Version]
\label{thm:audit_complexity}
Let auditor $\mathcal{A}$ make a claim $\mathcal{C}$ about SCX system state $s_t$ (e.g., ``$s_t$ was not tampered with at append operation $a_k$'').
Let the \textbf{evidence} for claim $\mathcal{C}$ on the CEC chain be $\mathcal{E} \subseteq \{s_0, \ldots, s_t\}$.
If there exists an effective verification algorithm $\mathcal{V}$ satisfying $\mathcal{V}(\mathcal{E}) = 1 \iff \mathcal{C}$ is true,
then the Kolmogorov complexity of the evidence set $\mathcal{E}$ satisfies:
\begin{equation}\label{eq:kolmogorov_audit}
\Kolm(\mathcal{E}) \geq \Kolm(\mathcal{C}) - \Kolm(\mathcal{V}) - O(1).
\end{equation}
\end{theorem}

\begin{proof}[Proof --- Basic AIT Inequality]
Let $p_{\mathcal{C}}$ be the shortest program generating $\mathcal{C}$ and $p_{\mathcal{V}}$ the code of the verification algorithm $\mathcal{V}$.
The standard AIT result is: if $\mathcal{C}$ can be effectively inferred from $\mathcal{E}$,
then $\Kolm(\mathcal{C}) \leq \Kolm(\mathcal{E}) + \Kolm(\text{inference program}) + O(1)$.
This is the definition of conditional Kolmogorov complexity: $\Kolm(\mathcal{C} \mid \mathcal{E}) \leq \Kolm(\mathcal{V})$.
From this, $\Kolm(\mathcal{E}) \geq \Kolm(\mathcal{C}) - \Kolm(\mathcal{C} \mid \mathcal{E})$ follows immediately.

\textbf{Conclusion:} ``The complexity of audit evidence is not less than the complexity of the claim'' is an AIT \textbf{tautology} ---
because the information of the claim $\mathcal{C}$ cannot arise from nothing.
The SCX context does not change this trivial fact.
\strictness{Tautology --- identity within the AIT framework, not SCX-specific}
\end{proof}

\begin{honestbox}
\textbf{The practical vacuity of the SCX-Kolmogorov claim:}

Equation (\ref{eq:kolmogorov_audit}) is literally correct --- but in practice it is \textbf{completely useless}, because:
\begin{itemize}
    \item $\Kolm(\cdot)$ is an \textbf{uncomputable function}. Any assertion that ``the audit requires $\Kolm(\mathcal{E}) \ge 10^6$'' cannot be verified.
    \item In real SCX systems, the length $|\mathcal{E}|$ (a computable naive upper bound) is far larger than $\Kolm(\mathcal{E})$ ---
    the raw data on the CEC chain is highly redundant.
    \item Hence using $|\mathcal{E}|$ in place of $\Kolm(\mathcal{E})$ yields a lower bound that is \textbf{at least $\sim \log |\mathcal{E}|$ times too loose}.
\end{itemize}

\textbf{The role of ``Kolmogorov complexity'' in the SCX literature is more rhetorical ---}
suggesting that audit evidence ``the less compressible, the harder to forge,'' but lacking an operational quantification method.
\textbf{The real information-theoretic constraints come from} Shannon entropy and mutual information (as discussed in fano\_scx.tex),
not the uncomputable Kolmogorov complexity.
\end{honestbox}

\begin{remark}[Minimum Description Length (MDL) as an Alternative]
If one abandons Kolmogorov complexity in favor of a \textbf{practically computable} framework: Minimum Description Length (MDL)
and Normalized Compression Distance (NCD) are more operationally meaningful for real SCX audits.
For example, one can use the zlib/gzip compression ratio of the CEC chain as a practical proxy for ``audit evidence redundancy.''
\textbf{The SCX theoretical framework should prioritize developing computable information measures rather than remaining in the uncomputable domain of AIT.}
\end{remark}

% ===========================================================================
\section{Optimal Hypothesis Set Length Lower Bound}
\label{sec:hypothesis_bound}
% ===========================================================================

\begin{definition}[Sufficiency of SCX Hypothesis Sets]
\label{def:sufficient_hypothesis}
Let the audit target be: from observation $Y$, decide whether the true state $s^*$ belongs to the ``compliant state set'' $\mathcal{H}_0 \subset \mathcal{S}$.
A hypothesis set $\mathcal{H} \subseteq \mathcal{S}$ is called \textbf{$\varepsilon$-sufficient}
if there exists a decision rule $\delta: \mathcal{Y} \to \{0,1\}$ such that:
\[
\sup_{s \in \mathcal{H}} \mathbb{P}_s(\delta = 1) \leq \alpha, \quad
\inf_{s \notin \mathcal{H}} \mathbb{P}_s(\delta = 1) \geq 1 - \beta,
\]
and $\alpha + \beta \leq \varepsilon$.
\end{definition}

\begin{theorem}[Information-Theoretic Lower Bound on Optimal Hypothesis Set Length]
\label{thm:hypothesis_lb}
In the SCX framework, if audit operations cover depth $d$ on the CEC chain (i.e., the earliest auditable state is $s_{t-d}$),
and the alphabet of chained append operations is $\mathcal{A}$, then under any fixed audit strategy family,
the \textbf{description length} of the optimal hypothesis set $\mathcal{H}^*$ achieving error $\varepsilon$ (denoted $|\mathcal{H}^*|_{\text{desc}}$) satisfies:
\begin{equation}\label{eq:hypothesis_lb}
|\mathcal{H}^*|_{\text{desc}} \geq d \cdot \log|\mathcal{A}| - \frac{h_2(\varepsilon)}{\varepsilon} + O(1).
\end{equation}
\end{theorem}

\begin{proof}[Proof]
\strictness{Sketch: combines Fano and Shannon source coding theorem; lower bound constants need optimization}
\begin{enumerate}
    \item A CEC chain of depth $d$ has $|\mathcal{A}|^d$ possible append sequences.
    Each sequence corresponds to a state --- hence $|\mathcal{A}|^d$ possible histories must be distinguished.

    \item If the description length of hypothesis set $\mathcal{H}$ is $< d\log|\mathcal{A}|$,
    then at least two different CEC histories map to the same hypothesis description (pigeonhole principle).
    These two histories are \textbf{indistinguishable} to the auditor.

    \item By Fano's inequality (\ref{eq:fano_weak}), among $M = |\mathcal{A}|^d$ equally likely histories,
    the error probability $P_e$ is lower bounded by the conditional entropy:
    \[
    P_e \geq 1 - \frac{I(\text{history}; \text{hypothesis description}) + \log 2}{\log |\mathcal{A}|^d}.
    \]
    Since $I \leq |\mathcal{H}^*|_{\text{desc}}$ (description length upper-bounds information),
    we obtain:
    \[
    |\mathcal{H}^*|_{\text{desc}} \geq (1 - P_e) \cdot d\log|\mathcal{A}| - \log 2.
    \]
    Setting $P_e = \varepsilon$ yields (\ref{eq:hypothesis_lb}).

    \item \textbf{Caveat:} The above uses the ``equally likely histories'' assumption --- which \textbf{does not necessarily hold} in SCX.
    Real CEC append sequence distributions are influenced by user behavior and are non-uniform.
    Under a non-uniform prior, \textbf{entropy} must replace $\log|\mathcal{A}|$:
    \[
    |\mathcal{H}^*|_{\text{desc}} \geq d \cdot \Ent(\text{append distribution}) \cdot (1 - \varepsilon) - \log 2.
    \]
\end{enumerate}
\end{proof}

\begin{limitbox}
\textbf{Tightness analysis of the lower bound:}

Equation (\ref{eq:hypothesis_lb}) provides an \textbf{$\Omega(d)$}-level lower bound, but:
\begin{itemize}
    \item Tightness is unknown: does there exist a hypothesis set description scheme that \textbf{achieves} this lower bound (within a constant factor)?
    This is equivalent to a ``hypothesis set source coding theorem'' for SCX --- currently \textbf{unsolved}.
    \item The constant factor $1 - \varepsilon$ is crude --- a more refined strong converse analysis
    can yield the improvement $1 - \varepsilon \to 1 - \frac{\varepsilon}{\log|\mathcal{A}|}$.
    \item \textbf{Practical significance:} For $d=10^6$, $|\mathcal{A}|=256$, $\varepsilon=0.01$,
    the lower bound is $\approx 7.92 \times 10^6$ bits ($\approx 1$ MB).
    This is \textbf{negligible} in modern storage --- the practical constraining power of this lower bound is \textbf{only effective for extremely depth-constrained audit scenarios}.
\end{itemize}
\end{limitbox}

\begin{theorem}[Improved Lower Bound Beyond the Uniform Prior --- Joint Fano-LeCam]
\label{thm:improved_lb}
Under a non-uniform prior $\pi$, the optimal hypothesis set length satisfies:
\begin{equation}\label{eq:improved_hyp_lb}
|\mathcal{H}^*|_{\text{desc}} \geq \sup_{Q \ll \pi} \left\{
d \cdot \mathbb{E}_{a \sim Q}[\Ent(P_{Y|s(a)})] - \KL(Q \| \pi)
\right\} \cdot (1 - \varepsilon) - \log 2,
\end{equation}
where $Q$ ranges over all forensic distributions absolutely continuous with respect to $\pi$, and $s(a)$ is the state corresponding to append sequence $a$.
\strictness{Rigorous: follows from the variational form of the Fano-LeCam joint bound, holds for any prior $\pi$}
\end{theorem}

% ===========================================================================
\section{SCX Generalization of Quantum Fisher Information}
\label{sec:quantum_fisher}
% ===========================================================================

\begin{assumption}[Existence Premise of Quantum SCX]\label{as:quantum_scx}
All discussion below assumes that the SCX framework has been consistently generalized to the quantum domain:
CEC states $\rho_s$ are density operators on a Hilbert space $\mathcal{H}_{\text{CEC}}$,
and append operations correspond to completely positive trace-preserving (CPTP) maps $\mathcal{E}_a: \rho_s \mapsto \rho_{s \circ a}$.
\textbf{Warning:} As of this writing, a complete axiomatic system for quantum SCX \textbf{has not yet been established}.
The following content is a \textbf{formal generalization} rather than an established theory.
\strictness{Conjecture: quantum SCX does not yet exist}
\end{assumption}

\begin{definition}[Classical Fisher Information --- Review]
\label{def:classical_fisher}
For a parametric family $\{P_\theta\}_{\theta \in \Theta}$ ($\Theta \subseteq \mathbb{R}$),
the Fisher information is:
\begin{equation}
\FisherClassical(\theta) = \mathbb{E}_{X \sim P_\theta}\left[
\left(\frac{\partial}{\partial \theta} \log p_\theta(X)\right)^2
\right].
\end{equation}
Cramér-Rao lower bound: $\Var(\hat{\theta}) \geq 1/\FisherClassical(\theta)$.
\end{definition}

\begin{definition}[Quantum Fisher Information --- Helstrom Form]
\label{def:quantum_fisher}
For a quantum state family $\{\rho_\theta\}_{\theta \in \Theta}$, the symmetric logarithmic derivative (SLD) $L_\theta$ satisfies:
\begin{equation}
\frac{\partial \rho_\theta}{\partial \theta} = \frac{1}{2}(L_\theta \rho_\theta + \rho_\theta L_\theta).
\end{equation}
The quantum Fisher information (SLD form) is:
\begin{equation}
\FisherQuantum(\theta) = \Tr(\rho_\theta L_\theta^2).
\end{equation}
Quantum Cramér-Rao bound: $\Var_\theta(\hat{\theta}) \geq 1/\FisherQuantum(\theta)$.
\strictness{Rigorous: fundamental theorem of quantum estimation theory}
\end{definition}

\begin{theorem}[SCX--Quantum Fisher Information --- Formal Generalization]
\label{thm:scx_qfi}
Under Assumption~\ref{as:quantum_scx}, consider a quantum state $\rho_d$ at depth $d$ on the CEC chain.
The auditor obtains classical output $Y$ via a POVM measurement $\{M_y\}_{y \in \mathcal{Y}}$,
and attempts to estimate the quantum Hellinger distance between the CEC state and the ``honest state'' $\rho_{\text{honest}}$.
Then the quantum limit of audit precision is:
\begin{equation}\label{eq:scx_qfi_bound}
\Var(\hat{\theta}) \geq \frac{1}{d \cdot \overline{\FisherQuantum}},
\end{equation}
where $\overline{\FisherQuantum} = \frac{1}{d}\sum_{i=1}^{d} \FisherQuantum_i$ is the average quantum Fisher information of each append operation along the chain (under the audit parametrization).
\end{theorem}

\begin{proof}[Derivation Sketch --- Quantum Cramér-Rao + Chain Additivity]
\strictness{Sketch: uses additivity of quantum Fisher information (rigorous for product states),
but requires correction for entangled CEC chains}
\begin{enumerate}
    \item If the quantum states along the CEC chain are \textbf{unentangled} across different append steps (i.e., $\rho_d = \bigotimes_{i=1}^{d} \rho_{(i)}$),
    then quantum Fisher information has strict additivity: $\FisherQuantum_{\text{total}} = \sum_i \FisherQuantum_i$.

    \item The Cramér-Rao bound directly gives: $\Var(\hat{\theta}) \geq 1/\FisherQuantum_{\text{total}} = 1/(d \cdot \overline{\FisherQuantum})$.

    \item If the CEC chain has \textbf{quantum entanglement} (append operations generate cross-step entanglement),
    then $\FisherQuantum_{\text{total}}$ can be \textbf{superadditive} ---
    $\FisherQuantum$ for a $d$-step entangled chain can reach $O(d^2)$ (Heisenberg scaling),
    in which case (\ref{eq:scx_qfi_bound}) becomes an \textbf{overly conservative lower bound}:
    \[
    \Var(\hat{\theta}) \geq \frac{1}{O(d^2) \cdot \overline{\FisherQuantum}} \quad (\text{entangled quantum CEC}).
    \]
\end{enumerate}
\end{proof}

\begin{honestbox}
\textbf{Fatal problems with the SCX generalization of quantum Fisher information:}

\begin{enumerate}[leftmargin=*]
    \item \textbf{Quantum SCX does not exist.} All discussion of ``quantum CEC chains'' is formal manipulation on a mathematical object that has not yet been defined.
    \item If the append operations of the CEC chain are \textbf{classical} (discrete symbol appends), then ``quantum CEC'' is a category error ---
    CEC is intrinsically a \textbf{logical/symbolic structure}, not a physical quantum state.
    \item The only way to map CEC to physical quantum states is to quantize the \textbf{storage medium} ---
    but this changes the storage layer, not the logical structure of CEC. A classical CEC chain on a quantum storage medium is still classical information.

    \item \textbf{A possible legitimate connection:} If SCX audits use \textbf{quantum sensing / quantum metrology} techniques
    to read the storage medium (e.g., using SQUIDs to read faint traces on magnetic media),
    then audit sensitivity is bounded by the quantum Fisher information --- but this is the limit of \textbf{quantum readout technology},
    not the limit of the SCX framework.
\end{enumerate}

\textbf{Recommendation:} Demote ``SCX generalization of quantum Fisher information'' to an \textbf{appendix} or ``future work'' ---
until quantum SCX is properly defined, this content can only exist as formal exploration and
\textbf{should not appear in the theorem list}.
\end{honestbox}

% ===========================================================================
\section{Synthesis of the Three Limits: The Panorama of SCX Information-Theoretic Boundaries}
\label{sec:synthesis}
% ===========================================================================

\begin{table}[h!]
\centering
\caption{\textbf{Nature and Status of the Three SCX Information-Theoretic Limits}}
\label{tab:three_limits}
\renewcommand{\arraystretch}{1.6}
\begin{tabular}{@{}p{2.8cm} p{3.5cm} p{3.5cm} p{3.5cm}@{}}
\toprule
& \textbf{Kolmogorov Complexity} & \textbf{Hypothesis Set Length LB} & \textbf{Quantum Fisher Info} \\
\midrule
\textbf{Core inequality} &
$\Kolm(\mathcal{E}) \geq \Kolm(\mathcal{C}) - \Kolm(\mathcal{V})$ &
$|\mathcal{H}^*|_{\text{desc}} \geq d \cdot \log|\mathcal{A}| \cdot (1-\varepsilon)$ &
$\Var(\hat{\theta}) \geq 1/(d \cdot \overline{\FisherQuantum})$ \\
\midrule
\textbf{Computable?} &
\textbf{Uncomputable} &
Computable (given model) &
Computable (given quantum state family) \\
\midrule
\textbf{SCX originality} &
\textbf{Zero} --- AIT restatement &
\textbf{Low} --- Fano-LeCam parametrized for SCX &
\textbf{Formal analogy} --- not yet defined \\
\midrule
\textbf{Rigor} &
$\Kolm(\cdot)$ uncomputable makes
\textbf{rigor meaningless} &
$\log$ factor precise,
constant factor $\sim 30\%$ &
\textbf{Premise does not exist} \\
\midrule
\textbf{Practical constraining power} &
\textbf{None} ---
uncomputable + extremely loose &
\textbf{Weak} ---
$\sim$ MB level, not a bottleneck &
\textbf{None} ---
quantum SCX does not exist \\
\midrule
\textbf{Recommended priority} &
\textbf{Low} ---
switch to MDL framework &
\textbf{Medium} ---
refine non-uniform prior &
\textbf{Suspend} ---
await quantum SCX definition \\
\bottomrule
\end{tabular}
\end{table}

\begin{limitbox}
\textbf{The Real Picture of SCX Information-Theoretic Limits:}

The information-theoretic kernel of SCX consists of \textbf{three concentric circles}:
\begin{enumerate}[leftmargin=*]
    \item \textbf{Inner circle (rigorous) --- Shannon information theory:} The Fano-LeCam-Fisher triad provides the \textbf{truly rigorous lower bounds} for SCX audits.
    These are applications of classical information theory parametrized for SCX --- not SCX-original but valid.

    \item \textbf{Middle circle (heuristic) --- AIT and thermodynamics:} Kolmogorov complexity + Landauer's principle provide \textbf{conceptual analogies and heuristic arguments},
    each with a fatal flaw in rigor or operability (uncomputable / orders-of-magnitude gap).

    \item \textbf{Outer circle (conjectural) --- quantum generalization:} Quantum Fisher information + quantum CEC belong to \textbf{mathematical conjectures not yet grounded}.
    Until quantum SCX is properly defined + quantum advantage is proven under that definition,
    this content does not belong to SCX's \textbf{established theory}.
\end{enumerate}

\textbf{Recommended self-discipline standards for SCX theory papers:}
\begin{itemize}
    \item Inner-circle results: may be called ``SCX Theorems'' or ``SCX lower bounds.''
    \item Middle-circle results: must be labeled ``heuristic'' or ``analogy'' and may not be listed as theorems.
    \item Outer-circle results: may only appear in ``Outlook / Discussion'' sections and may not be listed as theorems or main conclusions.
\end{itemize}
\end{limitbox}

% ===========================================================================
\section{Honest Punch Summary}
\label{sec:conclusion}
% ===========================================================================

\begin{honestbox}[title=Full-Paper Honest Punch Summary]
\begin{enumerate}[leftmargin=*]
    \item \textbf{Kolmogorov complexity hypothesis:}
    The ``Kolmogorov complexity lower bound for audit claims'' in the SCX literature is essentially a restatement of an AIT identity,
    and is operationally meaningless because $\Kolm(\cdot)$ is uncomputable.
    \textbf{It should be replaced by computable measures such as MDL/NCD/compression ratio.}

    \item \textbf{Optimal hypothesis set lower bound:}
    Theorem~\ref{thm:hypothesis_lb} is correct under a uniform prior,
    but the constraint it provides at practical scales ($\sim$ MB) \textbf{is not a bottleneck}.
    The real information bottleneck of SCX is the \textbf{linear growth of chain depth $d$}, not the description length of the hypothesis set.

    \item \textbf{Quantum Fisher information:}
    This section is a \textbf{formal void} until quantum SCX is defined.
    The current formal derivation (\ref{eq:scx_qfi_bound}) performs legitimate algebraic operations on a \textbf{non-existent mathematical object} ---
    analogous to computing the area formula of a ``round square.''
    \textbf{It is strongly recommended to suspend this direction until quantum SCX has a complete axiomatic definition.}

    \item \textbf{Overall assessment:}
    The rigorous part of SCX information-theoretic limits (Shannon/Fano/LeCam) is a restatement of classical information theory in SCX terms,
    the heuristic part (Kolmogorov/Landauer analogies) has educational value but is not theorem-grade,
    and the conjectural part (quantum Fisher) is a formal generalization without definition.
    \textbf{Honesty standards require SCX papers to clearly distinguish these three layers.}
\end{enumerate}
\end{honestbox}

% ===========================================================================
\bibliographystyle{plain}
\begin{thebibliography}{99}

\bibitem{li2008}
M.~Li and P.~M.~B.~Vitányi.
\textit{An Introduction to Kolmogorov Complexity and Its Applications, 3rd ed.}
Springer, 2008.

\bibitem{grunwald2007}
P.~D.~Grünwald.
\textit{The Minimum Description Length Principle.}
MIT Press, 2007.

\bibitem{cover2006}
T.~M.~Cover and J.~A.~Thomas.
\textit{Elements of Information Theory, 2nd ed.}
Wiley, 2006.

\bibitem{helstrom1976}
C.~W.~Helstrom.
\textit{Quantum Detection and Estimation Theory.}
Academic Press, 1976.

\bibitem{holevo1982}
A.~S.~Holevo.
\textit{Probabilistic and Statistical Aspects of Quantum Theory.}
North-Holland, 1982.

\bibitem{paris2009}
M.~G.~A.~Paris.
Quantum estimation for quantum technology.
\textit{Int. J. Quantum Inf.}, 7(supp01):125--137, 2009.

\bibitem{tsybakov2009}
A.~B.~Tsybakov.
\textit{Introduction to Nonparametric Estimation.}
Springer, 2009.

\bibitem{ciliberto2013}
S.~Ciliberto, A.~Imparato, A.~Naert, and M.~Tanase.
Heat flux and entropy produced by thermal fluctuations.
\textit{Phys. Rev. Lett.}, 110(18):180601, 2013.

\end{thebibliography}

\end{document}
